\documentstyle[% 


righttag% 


,11pt% 


,amssymb% 


,russian%               remove in Eng. 


,izvru%                 Set izven% in Eng. 


]{amsart} 


 


%  


\newcommand{\const}{\operatorname{const}} 


\newcommand{\sgrad}{\operatorname{sgrad}} 


\newcommand{\tts}[1]{} 


 


%\hoffset=-15mm%                  For Eng. 


 


\setcounter{page}{33} 


\god{2001} 


\nomer{3~(466)} %                        Remove in Eng. 


%\nomer{Vol.\,45, No.\,3}%               Set in Eng. 


%\str{\pageref{begin}--\pageref{end}}%  Set in Eng. 


\udk{514.762} 


 


\author{\it С.Г.\,Лейко} 


% NB:   ^^ remove in Eng. 


\title{О геодезическом потоке на сферическом  


касательном расслоении двумерного многообразия\\ с  


метрикой Сасаки} 


\thanks{Работа частично поддержана Международной Соросовской Программой  


поддержки образования в области точных наук (ISSEP), грант 


\No\,APU\,071057.} 


 


\date{} 


% \pagestyle{myheadings}%         Set in Eng. 


\pagestyle{plain} %              Remove in Eng. 


\begin{document} 


% \thispagestyle{plain}%          Set in Eng. 


\maketitle 


\label{begin} 


 


Как известно, гамильтоновы уравнения составляют один из важнейших 


классов дифференциальных уравнений. В частности, уравнения этого вида 


возникают в задаче нахождения геодезических кривых на римановых 


многообразиях. При этом среди всех гамильтоновых систем случай вполне 


интегрируемых представляется крайне редко [1]--[4]. В данной работе 


рассмотрен геодезический поток на сферическом касательном расслоении 


двумерного римановa многообразия с метрикой Сасаки и показано, что 


если базисное многообразие локально изометрично поверхности вращения, 


то соответствующая потоку гамильтонова система вполне интегрируема по 


Лиувиллю. Отсюда, как следствие, находятся траектории потока в 


квадратурах. 


 


Данное исследование возникло в связи с изучением автором вариационных 


задач для функционалов поворота кривых [5]--[9]. На основании 


результатов П.\,Надя [10] выяснилось, что базисные траектории потока 


(т.\,е. проекции траекторий потока на базу) являются 


изопериметрическими экстремалями поворота (первого порядка).


Найденный нами для 


этих экстремалей обобщенный интеграл Клеро оказался интегралом, которого


недоставало для полной интегрируемости геодезического потока 


на сферическом касательном расслоении с метрикой Сасаки [11], [12] в 


случае, когда базисное многообразие локально изометрично поверхности 


вращения. Отметим, что классический 


интеграл Клеро приводит к полной интегрируемости геодезического потока 


на поверхности вращения [1], [2]. Автор считает своим приятным долгом 


отметить, что гипотеза о такой роли обобщенного интеграла Клеро в 


теории геодезических потоков на расслоении была высказана А.Т.\,Фоменко 


и полностью подтвердилась. 


 


\section{Изопериметрические экстремали поворота\protect\\


на двумерных римановых многообразиях} 


 


В римановом пространстве $(M^n,g)$ рассмотрим функционал длины 


$l[\gamma]=\displaystyle\int^{t_0}_{t_1}(g_{i j}\Dot x^i\Dot 


x^j)^{1/2}dt$, где $\Dot x^i=\xi^i$ --- компоненты касательного вектора 


$\Dot\gamma$ 


некоторой параметризованной кривой $\gamma$, и функционал абсолютного 


поворота $\theta[\gamma]=\displaystyle\int^{l_0}_{l_1}k_g dl$. Здесь 


$l$ --- длина дуги на кривой $\gamma$, $k_g$ ---


первая кривизна Френе кривой $\gamma$, а в случае двумерного пространства


$k_g$ есть абсолютная геодезическая кривизна. 


 


Рассмотрим для функционала поворота изопериметрическую 


вариационную задачу с фиксированными концами: 


$$ 


\delta\theta=0,\ \; \gamma(t_0)=p_0,\ \; \gamma(t_1)=p_1, 


\ \; l_0=0,\ \; l_1=\wh l,\ \; l[\gamma]=\wh l=\const. 


$$ 


Применяя стандартный метод Эйлера--Лагранжа, убеждаемся в том,


что ее решения в двумерном пространстве удовлетворяют уравнению 


\begin{equation} 


k_g=c K, 


\end{equation} 


где $c$ --- постоянная, зависящая от фиксированной длины $\wh l$, $K$ 


--- гауссова 


кривизна пространства. Кроме того, в особом случае $K=0$, решением 


является всякая допустимая кривая (класса $C^4$ без точек 


спрямления). 


Кривые, удовлетворяющие уравнению (1), названы изопериметрическими 


экстремалями поворота пространства $(M^2,g)$ [6]. 


 


Отметим, что уравнение (1) ранее рассматривалось в исследованиях 


А.\,Пуанкаре в связи с изучением замкнутых геодезических кривых 


овальной 


поверхности, к которому сводилась астрономическая ``задача о трех 


телах'' 


([13], с.\,229). Нами установлены экстремальные свойства 


изопериметрических 


экстремалей поворота и получены их дифференциальные уравнения в 


нормальной форме [8]. В случае, если пространство $(M^2,g)$ есть


поверхность евклидова пространства, изопериметрическим экстремалям 


поворота дана механическая интерпретация [9]. На поверхностях вращения 


найден обобщенный интеграл Клеро 


\begin{equation} 


r\sin\omega=e c\sin\psi+c_1,\quad 


e=\pm1,\ \; c,c_1\text{ --- }\const, 


\end{equation} 


где $\omega$ --- угол между экстремалью и меридианом в их общей точке, 


$r$ --- 


расстояние от этой точки до оси вращения, $\psi$ --- угол, образованный 


касательной к меридиану с осью вращения [8]. 


 


\section{Геодезические кривые на сферическом касательном расслоении 


двумерного риманова многообразия с метрикой Сасаки} 


 


Метрика Сасаки впервые была рассмотрена на сферическом касательном 


расслоении единичных векторов $T_1M^n$, а затем П.\,Надь обобщил ее на 


$T_\rho M^n$ --- сферическое касательное расслоение векторов, квадрат 


длины которых равен постоянной $\rho$ [10]--[12]. 


 


Пусть $(M^2,g)$ --- произвольное риманово многообразие с метрикой 


$d l^2=g_{i j}d x^i\,d x^j$, $i,j,\dotsc=1,2$. 


Тогда на $T_\rho M^n$ индуцируется метрика 


\begin{equation} 


d t^2=g_{i j}d x^i\,d x^j+g_{i j}D y^i\,D y^j,\quad 


g_{i j}y^i y^j=\rho,\quad D y^k=d y^k+ 


\Gamma^k_{i j}y^i d x^j. 


\end{equation} 


 


Введем в координатной окрестности многообразия $M^2$ 


полугеодезические координаты $x^1$, $x^2$. Тогда 


$$ 


d l^2=(d x^1)^2+G(x^1,x^2)(d x^2)^2,\ \; 


K=-\frac{(\sqrt{G})_{11}}{\sqrt{G}}, 


\ \; \Gamma^1_{22}=-\frac{1}{2}G_1,\ \; \Gamma^2_{12}= 


\frac{G_1}{2G},\ \; \Gamma^2_{22}=\frac{G_2}{2G}, 


$$ 


а остальные символы Кристоффеля равны нулю. 


В качестве третьей координаты $x^3$ на $T_\rho M^2$ возьмем угол 


между касательным вектором $y^i$ и касательным вектором 


$\frac{\partial}{\partial x^1}$ к первой 


координатной линии в точке $(x^1,x^2)$. В этом случае 


$$ 


y^1=\sqrt{\rho}\,\cos x^3,\quad 


y^2=\frac{\sqrt{\rho}}{\sqrt{G}}\sin x^3. 


$$ 


В координатах $(x^1,x^2,x^3)$ метрика $d t^2$ на $T_\rho M^n$ в силу 


(3) 


приобретает вид 


$$ 


d t^2=g^*_{\alpha\beta}d x^\alpha\,d x^\beta, 


\ \;\alpha,\beta,\dotsc=1,2,3;\ \; (g^*_{\alpha\beta})= 


\begin{pmatrix} 


1 & 0 & 0 \\ 


0 & G+\rho(\sqrt{G})^2_1 & \rho(\sqrt{G})_1 \\ 


0 & \rho(\sqrt{G})_1 & \rho 


\end{pmatrix} 


. 


$$ 


Отсюда находим компоненты взаимного метрического тензора 


$g^*_{\alpha\beta}$ и 


символы Кристоффеля первого и второго рода 


\begin{gather*} 


(g^{*\alpha\beta})= 


\begin{pmatrix} 


1 & 0 & 0 \\ 


0 & G^{-1} & -G^{-1}(\sqrt{G})_1 \\ 


0 & -G^{-1}(\sqrt{G})_1 & \rho^{-1}+G^{-1}(\sqrt{G})^2_1 


\end{pmatrix} 


, \\ 


% 


\Gamma^*_{12,3}=\Gamma^*_{13,2}=-\Gamma^*_{23,1}= 


\frac{1}{2}\frac{\partial g^*_{23}}{\partial x^1}, 


\ \; \Gamma^*_{12,2}=-\Gamma^*_{22,1}=\frac{1}{2} 


\frac{\partial g^*_{22}}{\partial x^1}, \\


%


\Gamma^*_{22,2}=\frac{1}{2} 


\frac{\partial g^*_{22}}{\partial x^2}, 


\ \; \Gamma^*_{22,3}=\frac{\partial g^*_{23}}{\partial x^2} 


\ \; \text{(остальные --- нули)};\\ 


% 


\Gamma^{*1}_{\alpha\beta}=g^{*1\gamma}\Gamma^*_{\alpha\beta,\gamma}= 


\Gamma^*_{\alpha\beta,1}, 


\ \; \Gamma^{*2}_{\alpha\beta}=g^{*2\gamma} 


\Gamma^*_{\alpha\beta,\gamma}=G^{-1} 


\Gamma^*_{\alpha\beta,2}-G^{-1}(\sqrt{G})_1 


\Gamma^*_{\alpha\beta,3}, \\ 


% 


\Gamma^{*3}_{\alpha\beta}=g^{*3\gamma} 


\Gamma^*_{\alpha\beta,\gamma}=-G^{-1}(\sqrt{G})_1 


\Gamma^*_{\alpha\beta,2}+(\rho^{-1}+G^{-1}(\sqrt{G})^2_1) 


\Gamma^*_{\alpha\beta,3}. 


\end{gather*} 


 


Относительно введенных координат $x^\alpha$ элемент поворота $d\theta$ 


и геодезическая кривизна кривой $x^i(l)$ приобретают вид 


$$ 


d\theta=d x^3+(\sqrt{G})_1d x^2,\quad 


k_g=e\frac{d\theta}{d l}=e 


\bigg(\frac{d x^3}{d l}+(\sqrt{G})_1 


\frac{d x^2}{d l} \bigg),\quad e=\pm1. 


$$ 


Метрику $d t^2$ теперь можно представить в виде 


$$ 


d t^2=(d x^1)^2+(G+\rho(\sqrt{G})^2_1)(d x^2)^2+ 


2\rho(\sqrt{G})_1d x^2d x^3+\rho(d x^3)^2= 


d l^2+\rho d\theta^2. 


$$ 


 


Рассмотрим дифференциальные уравнения геодезических кривых 


$x^\alpha(t)$ на $T_\rho M^2$, отнесенных к натуральному параметру 


\begin{equation} 


\frac{d^2x^\gamma}{d t^2}+\Gamma^{*\gamma}_{\alpha\beta} 


\frac{d x^\alpha}{d t}\frac{d x^\beta}{d t}=0. 


\end{equation} 


Подставив сюда найденные значения символов Кристоффеля, получим 


\begin{gather} 


\frac{d^2x^1}{d t^2}+\Gamma^1_{i j} 


\frac{d x^i}{d t}\frac{d x^j}{d t}=\rho(\sqrt{G})_{11} 


\frac{d x^2}{d t}\frac{d\theta}{d t},\tag{4.1} \\ 


% 


\frac{d^2x^2}{d t^2}+\Gamma^2_{i j}\frac{d x^i}{d t} 


\frac{d x^j}{d t}=-\rho G^{-1}(\sqrt{G})_{11} 


\frac{d x^1}{d t}\frac{d\theta}{d t},\tag{4.2} \\ 


% 


\frac{d^2x^3}{d t^2}+G^{-1}\bigg[(G(\sqrt{G})_{11}- 


\rho(\sqrt{G})^2_1(\sqrt{G})_{11}-G_1(\sqrt{G})_1) 


\frac{d x^1}{d t}\frac{d x^2}{d t}+ \notag \\ 


% 


+(G(\sqrt{G})_{12}-\tfrac{1}{2}G_2(\sqrt{G})_1) 


\bigg(\frac{d x^2}{d t} \bigg)^2-\rho(\sqrt{G})_1 


(\sqrt{G})_{11}\frac{d x^1}{d t}\frac{d x^3}{d t}\bigg]=0.\tag{4.3} 


\end{gather} 


 


Из первых двух уравнений системы (4) вытекает, что вдоль геодезических 


кривых на $T_\rho M^2$ имеет место $g_{i j}\xi^i(t)\xi^j_1(t)=0$, 


$\xi^i_1=\nabla_t\xi^i$ и, следовательно, 


\begin{equation} 


h\equiv \tfrac{1}{2}g_{i j}\xi^i(t)\xi^j(t)=\const. 


\end{equation} 


Вычисляя геодезическую кривизну базисных кривых $x^i(t)$, 


вследствие (4), (5) получаем 


\begin{multline*} 


k^2_g(t)=\frac{\langle\xi,\xi \rangle 


\langle\xi_1,\xi_1 \rangle-\langle\xi,\xi_1 \rangle^2} 


{\langle\xi,\xi \rangle^3}=g_{i j}\xi^i_1(t)\xi^j_1(t) 


[g_{i j}\xi^i(t)\xi^j(t)]^{-2}= \\ 


% 


=(2h)^{-2}\bigg[\rho^2(\sqrt{G})^2_{11}\bigg( 


\frac{d x^2}{d t}\bigg)^2\bigg(\frac{d\theta}{d t} \bigg)^2+ 


\rho^2G^{-1}(\sqrt{G})^2_{11} 


\bigg(\frac{d x^1}{d t} \bigg)^2


\bigg(\frac{d\theta}{d t} \bigg)^2\bigg]= \\ 


% 


=\rho^2(2h)^{-2}G^{-1}(\sqrt{G})^2_{11}\bigg[G 


\bigg(\frac{d x^2}{d t} \bigg)^2+ 


\bigg(\frac{d x^1}{d t} \bigg)^2\bigg] 


\bigg(\frac{d\theta}{d t} \bigg)^2= 


\rho^2K^2(2h)^{-1}\bigg(\frac{d\theta}{d t} \bigg)^2. 


\end{multline*} 


Поскольку параметр $t$ выбран натуральным, то 


$$ 


g^*_{\alpha\beta}\xi^\alpha(t)\xi^\beta(t)= 


g_{i j}\xi^i(t)\xi^j(t)+\rho 


\bigg(\frac{d\theta}{d t} \bigg)^2=1, 


$$ 


а значит, вдоль геодезических 


\begin{gather} 


2h+\rho a^2=1, \\ 


a\equiv\frac{d\theta}{d t}=\frac{d x^3}{d t}+ 


(\sqrt{G})_1\frac{d x^2}{d t}=\const. 


\end{gather} 


Следовательно, 


\begin{equation} 


k_g=\frac{e\rho a}{\sqrt{2h}}K. 


\end{equation} 


Нетрудно убедиться в том, что вследствие (7) уравнение (4.3) 


вытекает из уравнений (4.1), (4.2). 


 


Равенство (7) дает промежуточный интеграл геодезических и означает, 


что касательный вектор $\frac{d x^i}{d t}$ вдоль базисной кривой 


$x^i(t)$ совершает 


простое винтовое движение: $\theta=at+a_0$, $a,a_0$ --- $\const$. В 


свою очередь, 


равенство (8) показывает, что базисная кривая является 


изопериметрической 


экстремалью поворота пространства $(M^2,g)$ с изопериметрической 


постоянной $c=\frac{e\rho a}{\sqrt{2h}}$.


 


Таким образом, результату из [10] можно придать следующую 


формулировку.


 


\begin{thm} 


Для того чтобы кривая $x^\alpha(t)$ была геодезической в 


сферическом касательном расслоении $T_\rho M^2$ с метрикой Сасаки 


$d{}t^2$, 


необходимо и $($при $K\ne0)$ достаточно, чтобы базисная кривая $x^i(t)$ 


была 


изопериметрической экстремалью поворота и ее касательный вектор 


$\frac{d x^i}{d t}$ 


совершал вдоль нее простое винтовое движение. 


\end{thm} 


 


\section{Геодезический поток на $T_\rho M^2$} 


 


Рассмотрим кокасательное расслоение $T^*(T_\rho M^2)$ с локальными 


координатами $x^\alpha$, $p_\alpha$ и естественной канонической 


симплектической 


структурой $\omega=d p_\alpha\wedge d x^\alpha$. Возьмем функцию 


Гамильтона $H(x,p)=\frac{1}{2}g^{*\alpha\beta}p_\alpha p_\beta$ 


и отвечающий ей гамильтонов поток $\Dot x=\sgrad H$ относительно 


симплектической структуры $\omega$ на $T^*(T_\rho M^2)$. Поскольку $H$ 


--- 


первый интеграл этого потока, то единичное кокасательное расслоение


$T^*_1(T_\rho M^2)=\{x^*\in T^*(T_\rho M^2):\|p\|=1\}$ инвариантнo 


относительно потока 


$\sgrad H$. Ограничение этого потока на $T^*_1(T_\rho M^2)$ будет 


геодезическим 


потоком на $T_\rho M^2$. При естественном изоморфизме $T^*(T_\rho 


M^2)\to T(T_\rho M^2)$ 


траектории геодезического потока $\sgrad H$ переходят в кривые, 


составленные из касательных векторов в $T_\rho M^2$. 


 


Пусть $\Phi^t$ --- локальная однопараметрическая группа преобразований, 


порожденная потоком $\sgrad H$. Отдельные преобразования из $\Phi^t$ 


переводят пару $(x(0),p(0))$ в пару $(x(t),p(t))=\Phi^t(x(0),p(0))$, 


где для 


получения $x(t)$ следует провести геодезическую через точку $x(0)$ в 


направлении ковектора $p(0)$, и тогда $x(t)$ отстоит от $x(0)$ на 


расстоянии 


$t$ вдоль этой геодезической, а ковектор $p(t)$ касается этой 


геодезической 


в $x(t)$ и направлен так же, как и $p(0)$. Таким образом, при получении 


$x(t)=(x^\alpha(t))$ на базе следует проводить изопериметрическую 


экстремаль 


поворота через исходную точку $x^i(0)$. С этой целью интегрируем 


систему 


(4.1), (4.2) с учетом промежуточного интеграла (7) и постоянную $a$ 


определим начальными данными 


$$ 


a=\frac{d x^3(0)}{d t}+(\sqrt{G})_1\frac{d x^2(0)}{d t}. 


$$ 


Наконец, для определения компоненты траектории $x^3(t)$ 


интегрируем уравнение (7) при вышеуказанной постоянной $a$. 


 


Как известно, гамильтониан $H(x,p)$ является основным первым 


интегралом геодезического потока. Из (7) имеем еще один первый интеграл 


$$ 


a=p_\alpha(g^{*3\alpha}+(\sqrt{G})_1 g^{*2\alpha})= 


\rho^{-1}p_3. 


$$ 


Аналогичным образом из (5) получаем первый интеграл 


$h=\frac{1}{2}g_{i j}g^{*i\alpha}g^{*j\beta}p_\alpha p_\beta$, 


однако в силу (6) интегралы $H$, $a$, $h$ зависимы: $2H=2h+\rho a^2$. 


 


В случае, когда многообразие $(M^2,g)$ локально изометрично 


поверхности вращения, возможно указать дополнительный интеграл потока, 


который вытекает из найденного нами обобщенного интеграла Клеро. 


Действительно, пусть многообразие $(M^2,g)$ локально изометрично 


поверхности вращения с меридианом $f(r)$, где $r$ --- расстояние до оси 


вращения. Тогда $d l^2=F^2d r^2+r^2(d x^2)^2$, $F=\sqrt{1+f'{}^2}$, 


где $x^2$ --- долгота точки. Вводя новую координату $x^1:d 


x^1=F\,d{}r$, 


получим метрику $d l^2$ в виде 


$$ 


d l^2=(d x^1)^2+r^2(d x^2)^2, 


$$ 


т.\,е. координаты $x^1$, $x^2$ полугеодезические и $r(x^1)=\sqrt{G}$. 


Нетрудно 


убедиться в том, что $\sin\psi=F^{-1}=(\sqrt{G})_1$. Поскольку 


для базисной траектории изопериметрическая постоянная


$c=\frac{e\rho a}{\sqrt{2h}}=\frac{e p_3}{\sqrt{2h}}$, то 


обобщенный интеграл Клеро (2) приобретает вид 


$$ 


\varkappa=\sqrt{G}\,\sin x^3+\frac{p_3}{\sqrt{2h}}. 


$$ 


Учитывая, что $\sin x^3=\sqrt{\frac{G}{2h}}\frac{d x^2}{d t}$, получим 


$$ 


\varkappa=\frac{G}{\sqrt{2h}}g^{*2\alpha}p_\alpha+ 


\frac{p_3}{\sqrt{2h}}(\sqrt{G})_1=\frac{p_2}{\sqrt{2h}}. 


$$ 


Отсюда вытекает, что в данном случае вторая компонента $p_2$ импульса 


также является интегралом геодезического потока. 


 


Рассмотрим скобку Пуассона канонической симплектической структуры 


$$ 


\{F_1,F_2\}=\sum_\alpha 


\frac{\partial F_1}{\partial p_\alpha} 


\frac{\partial F_2}{\partial x^\alpha}- 


\frac{\partial F_2}{\partial p_\alpha} 


\frac{\partial F_1}{\partial x^\alpha}. 


$$ 


Так как в рассматриваемом случае гамильтониан $H$ не зависит сразу от 


двух 


переменных $x^2$, $x^3$, то нетрудно проверить, что интегралы $H$, 


$p_2$, $p_3$ 


находятся в инволюции, т.\,е. 


$$ 


\{H,p_2\}=\{H,p_3\}=\{p_2,p_3\}=0. 


$$ 


Очевидно, указанные три интеграла независимы и разрешимы 


относительно импульсов $p_1$, $p_2$, $p_3$. Следовательно, выполнена 


теорема 


Лиувилля [1], [2] и функции $H$, $p_2$, $p_3$ образуют полное 


инволютивное 


семейство интегралов гамильтоновых уравнений 


$$ 


\frac{d x^\alpha}{d t}=\frac{\partial H}{\partial p_\alpha},\quad 


\frac{d p_\alpha}{d t}=-\frac{\partial H}{\partial x^\alpha}. 


$$ 


Тем самым имеет место 


 


\begin{thm} 


Если риманово многообразие $(M^2,g)$ локально 


изометрично поверхности вращения, то геодезический поток сферического 


касательного расслоения $T_\rho M^2$ с метрикой Сасаки является вполне 


интегрируемым. 


\end{thm} 


 


Как вытекает из результатов работы [8], соответствующая квадратура 


для базисных траекторий (не являющихся параллелями) имеет в координатах 


радиус--долгота вид 


$$ 


x^2=\int^r_{r_0} 


\frac{e c+c_1F}{r(r^2-(e c F^{-1}+c_1)^2)^{1/2}}d r+x^2_0. 


$$ 


При $c=0$ эта квадратура совпадает с известной квадратурой для 


геодезических кривых на поверхности вращения [14]. 


 


\begin{thebibliography}{99} 


 


\bibitem{1} 


Фоменкo А.Т. {\em Симплектическая геометрия. Методы и приложения}. -- 


М.: Изд-во МГУ, 1988. -- 413 с. 


\bibitem{2} 


Трофимов В.В., Фоменко А.Т. {\em Алгебра и геометрия интегрируемых 


гамильтоновых дифференциальных уравнений}. 


-- М.: Факториал, 1995. -- 448 с. 


\bibitem{3} 


{\em Труды семинара по векторному и тензорному анализу с их 


приложениями к геометрии, механике и физике}. -- М.: Изд-во МГУ, 1993. 


-- Вып.\, 25. -- Ч.\,1. -- 127 с. 


\bibitem{4} 


{\em Труды семинара по векторному и тензорному анализу с их 


приложениями к геометрии, механике и физике}. -- М.: Изд-во МГУ, 1993. 


-- Вып.\,25. -- Ч.\,2. -- 151 с. 


\bibitem{5} 


Лейко С.Г. {\em Вариационные задачи для функционалов поворота и 


спин-отображения псевдоримановых пространств} // 


Изв. вузов. Математика. -- 1990. -- \No\,10. -- С.\,9--17. 


\bibitem{6} 


Лейко С.Г. {\em Поворотные диффеоморфизмы на поверхностях евклидова 


пространства} // Матем. заметки. -- 1990. -- Т.\,47. -- \No\,3. 


-- С.\,52--57. 


\bibitem{7} 


Лейко С.Г. {\em Экстремали функционалов поворота кривых псевдориманова 


пространства и траектории спин-частиц в гравитационных полях} // 


Докл. РАН. -- 1992. -- Т.\,325. -- \No\,4. -- С.\,659--664. 


\bibitem{8} 


Лейко С.Г. {\em Изопериметрические экстремали поворота на поверхностях 


в евклидовом пространстве $E^3$} // 


Изв. вузов. Математика. -- 1996. -- \No\,6. -- С.\,25--32. 


\bibitem{9} 


Лейко С.Г. {\em Механическая интерпретация изопериметрических 


экстремалей поворота на поверхностях} // Вicник Одеськ. держ. ун-ту. -- 


1999. -- Т.\,4. -- \No\,4. -- С.\,102--105. 


\bibitem{10} 


Nagy P. {\em On the tangent sphere bundle of Riemannian $2$-manifold} // 


T\^ohoku Mat. J. -- 1977. -- V.\,29. -- P.\,203--208. 


\bibitem{11} 


Sasaki S. {\em On the differential geometry of tangent bundles of 


Riemannian manifolds}. 1 // T\^ohoku Mat. J. -- 1958. -- V.\,10. -- 


P.\,338--354. 


\bibitem{12} 


Sasaki S. {\em On the differential geometry of tangent bundles of 


Riemannian manifolds}. 2 // T\^ohoku Mat. J. -- 1962. -- V.\,14. -- 


P.\,146--155. 


\bibitem{13} 


Бляшке В. {\em Дифференциальная геометрия}. -- М.--Л.: ОНТИ--НКТП 


СССР, 1935. -- 332 с. 


\bibitem{14} 


Каган В.Ф. {\em Основы теории поверхностей в тензорном изложении. 


Аппарат исследования, общие основания теории и внутренняя 


геометрия поверхности}. Ч.\,1. -- М.--Л.: ГИТТЛ, 1947. -- 512 с. 


 


\end{thebibliography} 


 


\vskip 2cm 


 


\post{\it Одесский государственный университет}{\it Поступила} 


\post{}{27.10.1998} 


 


\label{end} 


\end{document} 


